\chapter{Au-delà du regard}
\label{chap:au-dela-du-regard}

Lorsque Dürzo rejoignit quelque temps la caravane, Kael remarqua rapidement
qu'il ne ressemblait à aucun des voyageurs qu'il avait rencontrés jusque-là.

Ce n'était pas vraiment son apparence. L'homme ne portait rien qui permette de
comprendre immédiatement qui il était ou ce qu'il faisait. Il parlait peu,
restait généralement à l'écart et ne semblait chercher la compagnie de
personne. Pourtant, quelque chose chez lui mettait Kael mal à l'aise.

Pas exactement de la peur.

Plutôt cette sensation instinctive qui lui disait de garder ses distances.

Kael était pourtant devenu curieux de presque tout le monde. Lorsqu'un voyageur
possédait une compétence qu'il ne connaissait pas, il finissait généralement
par l'approcher, observer ce qu'il faisait et poser suffisamment de questions
pour découvrir s'il pouvait apprendre quelque chose.

Avec Dürzo, il n'en avait aucune envie. Ou plutôt, il en avait envie sans
parvenir à se décider à le faire.

Alors il l'observa de loin.

Ce fut ainsi qu'il remarqua quelque chose d'étrange : il avait du mal à le
suivre du regard.

Kael avait appris à observer. Cal lui avait montré comment surveiller quelqu'un
sans fixer constamment ses mouvements, comment utiliser les ombres, les sons et
tout ce qui se trouvait autour de lui. Les années passées sur les routes avaient
encore renforcé ces habitudes.

Pourtant, avec Dürzo, cela ne fonctionnait pas.

Kael pouvait savoir parfaitement où il se trouvait, détourner les yeux quelques
secondes, puis constater qu'il n'était plus là. Au début, il pensa simplement
avoir été distrait.

Puis cela recommença.

Et encore.

Il essaya alors de faire davantage attention, sans jamais s'approcher ni
chercher réellement à suivre l'homme. Il se contentait de l'observer depuis
l'autre côté du campement, parfois pendant plusieurs minutes.

Dürzo finissait toujours par lui échapper.

Un après-midi, Kael l'avait aperçu près de plusieurs chariots et décida de ne
plus le quitter des yeux. L'homme discutait avec un voyageur. Rien ne semblait
pouvoir le masquer et Kael voyait parfaitement les passages qu'il aurait pu
emprunter.

Quelqu'un passa devant lui avec une caisse.

Une seconde à peine.

Lorsque Kael regarda de nouveau, Dürzo n'était plus là. Il chercha autour des
chariots, puis derrière eux.

Rien.

— Cherches-tu quelqu'un ?

Kael se retourna.

Dürzo se tenait quelques mètres derrière lui.

Le garçon resta figé. Toutes les questions qu'il aurait normalement posées à
n'importe quel autre voyageur disparurent aussitôt.

Dürzo le regarda. Kael ouvrit la bouche.

Il bredouilla quelque chose qui ressemblait vaguement à un « non ».

L'homme soutint son regard encore un instant, puis repartit. Kael attendit
qu'il se soit éloigné avant de respirer normalement.

À partir de ce jour, il continua de l'observer.

Mais toujours de loin.

\scenebreak

Bartiméus, lui aussi, avait remarqué Dürzo.

Le chat accordait rarement beaucoup d'importance aux nouveaux voyageurs. Il
pouvait décider qu'il appréciait quelqu'un et s'installer sur ses affaires sans
autorisation, en ignorer totalement un autre ou, plus rarement, éviter une
personne dont la présence semblait lui déplaire.

Avec Dürzo, son comportement était différent. Bartiméus ne l'évitait pas, mais
ne cherchait pas davantage à obtenir son attention.

La première fois qu'il s'approcha de lui, le chat s'arrêta à quelques pas et
l'observa longuement. Dürzo baissa les yeux vers lui. Pendant plusieurs
secondes, aucun des deux ne bougea.

Puis Bartiméus s'assit.

Kael, qui assistait à la scène depuis un peu plus loin, attendit presque de
voir lequel des deux allait se lasser le premier.

Ce fut Dürzo.

Il détourna simplement le regard et reprit ce qu'il faisait. Bartiméus resta
encore quelques instants à l'observer avant de repartir à son tour.

À partir de ce jour, le chat conserva envers lui cette attitude inhabituelle.
Il pouvait passer à proximité de Dürzo sans la moindre crainte et ne refusait
pas de s'installer dans le même espace que lui. Pourtant, contrairement à son
habitude, il ne venait jamais se frotter contre ses jambes, ne cherchait pas à
grimper sur ses affaires et ne semblait pas particulièrement désireux de le
déranger.

Ce n'était pas de la peur.

C'était presque de la prudence.

\scenebreak

Kamatsu finit par remarquer l'intérêt que Kael portait à Dürzo.

— Depuis quand refuses-tu de poser une question à quelqu'un ?

— Je ne refuse pas.

— Alors va lui demander.

Kael resta silencieux.

Kamatsu sourit. Il connaissait depuis longtemps une manière particulièrement
efficace de faire réagir son ami.

— De toute façon, je suis sûr que tu n'es pas capable d'aller lui demander.

Kael tourna enfin la tête vers lui.

Pendant un instant, Kamatsu s'attendit à voir apparaître cette expression qu'il
connaissait si bien, celle qui précédait généralement un « cap » et une décision
plus ou moins raisonnable.

Mais Kael regarda simplement Dürzo.

Puis il secoua la tête.

— Non.

Le sourire de Kamatsu disparut peu à peu.

— Ah.

Kael fronça les sourcils.

— Quoi ?

— Rien.

Kamatsu observa Dürzo à son tour, puis revint vers son ami.

— C'est juste que... d'habitude, ça marche.

Kael ne répondit pas.

— T'intimide-t-il vraiment à ce point ?

Kael prit quelques secondes avant de répondre.

— Je ne sais pas.

Ce n'était pas de la peur. Du moins, il ne le pensait pas. Dürzo ne l'avait
jamais menacé et ne lui avait rien fait. Pourtant, chaque fois que Kael
envisageait de s'approcher de lui, quelque chose lui soufflait qu'il valait
mieux rester où il était.

Kamatsu sembla comprendre qu'il n'obtiendrait rien de plus.

— D'accord.

Ils restèrent encore un moment à regarder l'autre côté du campement. Quelqu'un
passa entre eux et Dürzo.

Lorsque le passage fut de nouveau dégagé, l'homme n'était plus là.

Kael soupira.

— Merde.

Kamatsu éclata de rire.

— Ça, par contre, ça n'a pas changé.

\scenebreak

Les jours suivants ne lui apportèrent aucune réponse.

Kael continua d'observer Dürzo lorsqu'il en avait l'occasion, mais il ne tenta
plus jamais de s'approcher suffisamment pour risquer d'être surpris comme la
première fois. Il cherchait simplement à comprendre ses déplacements, à repérer
ce qu'il faisait différemment des autres et surtout à découvrir comment un
homme pouvait être parfaitement visible un instant et devenir si difficile à
retrouver le suivant.

Il n'y parvint pas.

Parfois, il pensait avoir compris quelque chose : une manière de se placer, de
profiter du passage d'un autre voyageur ou de choisir un chemin que personne ne
regardait vraiment. Puis Dürzo faisait autre chose et toutes ses certitudes
disparaissaient.

Bartiméus continuait également à l'observer de temps en temps. Ce détail
agaçait presque Kael. Le chat ne paraissait pas davantage comprendre Dürzo que
lui, mais il semblait au moins savoir où regarder.

Quelques jours plus tard, Dürzo avait quitté la caravane.

Kael ne le vit pas partir.

Il s'en rendit compte lorsqu'il chercha machinalement sa silhouette parmi les
voyageurs et comprit qu'elle n'était plus là. Il demanda simplement à
Kamatsu :

— Dürzo est-il parti ?

— Ce matin.

Kael regarda la route derrière eux.

Il aurait voulu savoir comment il faisait. Il aurait voulu comprendre ce qu'il
avait manqué chaque fois que son regard avait cessé de suffire.

Mais il n'avait jamais osé le lui demander.

Et, d'une certaine manière, Kael préférait presque qu'il en soit ainsi.

Cal lui avait appris à observer. Ellundril lui avait montré ce que pouvaient
devenir le mouvement et la maîtrise lorsqu'ils semblaient parfaitement
naturels.

Dürzo, lui, ne lui avait rien appris.

Il lui avait simplement prouvé qu'il existait encore des choses qu'il ne savait
même pas comment regarder.